\documentclass{amsart}
% For dealing with references we use the comment environment
\usepackage{verbatim}
\newenvironment{reference}{\comment}{\endcomment}
%\newenvironment{reference}{}{}
\newenvironment{slogan}{\comment}{\endcomment}
\newenvironment{history}{\comment}{\endcomment}

% For commutative diagrams we use Xy-pic
\usepackage[all]{xy}

% We use 2cell for 2-commutative diagrams.
\xyoption{2cell}
\UseAllTwocells

% We use multicol for the list of chapters between chapters
\usepackage{multicol}

% This is generally recommended for better output
\usepackage{lmodern}
\usepackage[T1]{fontenc}

% For cross-file-references

% Package for hypertext links:
\usepackage{hyperref}

% For any local file, say "hello.tex" you want to link to please

% Theorem environments.
%
\theoremstyle{plain}
\newtheorem{theorem}[subsection]{Theorem}
\newtheorem{proposition}[subsection]{Proposition}
\newtheorem{lemma}[subsection]{Lemma}

\theoremstyle{definition}
\newtheorem{definition}[subsection]{Definition}
\newtheorem{example}[subsection]{Example}
\newtheorem{exercise}[subsection]{Exercise}
\newtheorem{situation}[subsection]{Situation}

\theoremstyle{remark}
\newtheorem{remark}[subsection]{Remark}
\newtheorem{remarks}[subsection]{Remarks}

\numberwithin{equation}{subsection}

% Macros
%
\def\lim{\mathop{\mathrm{lim}}\nolimits}
\def\colim{\mathop{\mathrm{colim}}\nolimits}
\def\Spec{\mathop{\mathrm{Spec}}}
\def\Hom{\mathop{\mathrm{Hom}}\nolimits}
\def\Ext{\mathop{\mathrm{Ext}}\nolimits}
\def\SheafHom{\mathop{\mathcal{H}\!\mathit{om}}\nolimits}
\def\SheafExt{\mathop{\mathcal{E}\!\mathit{xt}}\nolimits}
\def\Sch{\mathit{Sch}}
\def\Mor{\mathop{\mathrm{Mor}}\nolimits}
\def\Ob{\mathop{\mathrm{Ob}}\nolimits}
\def\Sh{\mathop{\mathit{Sh}}\nolimits}
\def\NL{\mathop{N\!L}\nolimits}
\def\CH{\mathop{\mathrm{CH}}\nolimits}
\def\proetale{{pro\text{-}\acute{e}tale}}
\def\etale{{\acute{e}tale}}
\def\QCoh{\mathit{QCoh}}
\def\Ker{\mathop{\mathrm{Ker}}}
\def\Im{\mathop{\mathrm{Im}}}
\def\Coker{\mathop{\mathrm{Coker}}}
\def\Coim{\mathop{\mathrm{Coim}}}

% Boxtimes
%
\DeclareMathSymbol{\boxtimes}{\mathbin}{AMSa}{"02}

%
% Macros for moduli stacks/spaces
%
\def\QCohstack{\mathcal{QC}\!\mathit{oh}}
\def\Cohstack{\mathcal{C}\!\mathit{oh}}
\def\Spacesstack{\mathcal{S}\!\mathit{paces}}
\def\Quotfunctor{\mathrm{Quot}}
\def\Hilbfunctor{\mathrm{Hilb}}
\def\Curvesstack{\mathcal{C}\!\mathit{urves}}
\def\Polarizedstack{\mathcal{P}\!\mathit{olarized}}
\def\Complexesstack{\mathcal{C}\!\mathit{omplexes}}
% \Pic is the operator that assigns to X its picard group, usage \Pic(X)
% \Picardstack_{X/B} denotes the Picard stack of X over B
% \Picardfunctor_{X/B} denotes the Picard functor of X over B
\def\Pic{\mathop{\mathrm{Pic}}\nolimits}
\def\Picardstack{\mathcal{P}\!\mathit{ic}}
\def\Picardfunctor{\mathrm{Pic}}
\def\Deformationcategory{\mathcal{D}\!\mathit{ef}}

\setlength{\emergencystretch}{3em}
\begin{document}
\title[Stacks Project, Tag 0APW]{723 --- Countably indexed affine completions with countably many defining functions are adic-star}
\maketitle
\noindent Copyright (C) 2005--2025 Johan de Jong. Stacks Project authors. Source: \href{https://stacks.math.columbia.edu/tag/0APW}{Tag 0APW}.

\noindent Editorial difficulty note (not author text). Difficulty H3: The proof constructs an uncountable-ordinal counterexample and proves the positive direction by a closure theorem for finitely generated ideals plus a diagonal exponent contradiction. The combined topological-algebra and cardinal construction is specialized research-level formal geometry..

\noindent Editorial provenance note (not author text). History: excerpt from revision \texttt{a0444\allowbreak 6e57e\allowbreak c1fbc\allowbreak 252a8\allowbreak 71afc\allowbreak ec775\allowbreak 2fb28\allowbreak 07b14}, formal-spaces.tex, lines 3094--3229. Reference presentation was adapted; original-author context and core support blocks were retained or added. The mathematical wording is unchanged. Licensed under GNU FDL 1.2 or later; no invariant sections or cover texts. See ../licenses/COPYING. Context blocks reproduce author definitions and supporting results; they are not additional counted problems.

\begin{definition}
\label{definition-types-affine-formal-algebraic-space}
Let $S$ be a scheme. Let $X$ be an affine formal algebraic space over $S$.
We say $X$ is {\it McQuillan} if $X$ satisfies the equivalent conditions
of Lemma \href{https://stacks.math.columbia.edu/tag/0AIC}{lemma mcquillan affine formal algebraic space (Tag 0AIC)}. Let $A$
be the weakly admissible topological ring associated to $X$. We say
\begin{enumerate}
\item $X$ is {\it classical} if $X$ is McQuillan and $A$ is admissible
(More on Algebra, Definition \ref{more-algebra-definition-topological-ring}),
\item $X$ is {\it weakly adic} if $X$ is McQuillan and $A$ is weakly adic
(Definition \ref{definition-weakly-adic}),
\item $X$ is {\it adic} if $X$ is McQuillan and $A$ is adic
(More on Algebra, Definition \ref{more-algebra-definition-topological-ring}),
\item $X$ is {\it adic*} if $X$ is McQuillan, $A$ is adic, and $A$
has a finitely generated ideal of definition, and
\item $X$ is {\it Noetherian} if $X$ is McQuillan and $A$ is
both Noetherian and adic.
\end{enumerate}
\end{definition}

\begin{definition}[Topological rings]
\label{more-algebra-definition-topological-ring}
\begin{reference}
\href{https://stacks.math.columbia.edu/bibliography}{Bibliography: EGA1 (Chapter 0, Sections 7.1 and 7.2)}
\end{reference}
Let $R$ be a ring and let $M$ be an $R$-module.
\begin{enumerate}
\item We say $R$ is a {\it topological ring} if $R$ is endowed with a topology
such that both addition and multiplication are continuous as maps
$R \times R \to R$ where $R \times R$ has the product topology.
In this case we say $M$ is a {\it topological module} if $M$ is endowed
with a topology such that addition $M \times M \to M$ and
scalar multiplication $R \times M \to M$ are continuous.
\item A {\it homomorphism of topological modules} is just a continuous
$R$-module map. A {\it homomorphism of topological rings} is a
ring homomorphism which is continuous for the given topologies.
\item We say $M$ is {\it linearly topologized} if $0$ has a fundamental
system of neighbourhoods consisting of submodules. We say $R$ is
{\it linearly topologized} if $0$ has a fundamental system of neighbourhoods
consisting of ideals.
\item If $R$ is linearly topologized, we say that $I \subset R$ is an
{\it ideal of definition} if $I$ is open and if every neighbourhood
of $0$ contains $I^n$ for some $n$.
\item If $R$ is linearly topologized, we say that $R$ is {\it pre-admissible}
if $R$ has an ideal of definition.
\item If $R$ is linearly topologized, we say that $R$ is {\it admissible} if
it is pre-admissible and
complete\footnote{By our conventions this includes separated.}.
\item If $R$ is linearly topologized, we say that $R$ is {\it pre-adic} if
there exists an ideal of definition $I$ such that $\{I^n\}_{n \geq 0}$
forms a fundamental system of neighbourhoods of $0$.
\item If $R$ is linearly topologized, we say that $R$ is {\it adic} if
$R$ is pre-adic and complete.
\end{enumerate}
Note that a (pre)adic topological ring is the same thing as a (pre)admissible
topological ring which has an ideal of definition $I$ such that $I^n$ is
open for all $n \geq 1$.
\end{definition}

\begin{definition}
\label{definition-weakly-adic}
\begin{reference}
\href{https://stacks.math.columbia.edu/bibliography}{Bibliography: Gabber-Ramero (Definition 8.3.8)}
\end{reference}
Let $A$ be a linearly topologized ring.
\begin{enumerate}
\item We say $A$ is {\it weakly pre-adic}\footnote{In \href{https://stacks.math.columbia.edu/bibliography}{Bibliography: Gabber-Ramero} the
authors say $A$ is {\it $c$-adic}.} if there exists an ideal
$I \subset A$ such that the closure of $I^n$ is open for all $n \geq 0$
and these closures form a fundamental system of open ideals.
\item We say $A$ is {\it weakly adic} if $A$ is weakly pre-adic
and complete\footnote{By our conventions this includes separated.}.
\end{enumerate}
\end{definition}

\begin{lemma}
\label{lemma-completion-countably-indexed}
\begin{reference}
Email by Ofer Gabber of September 11, 2014.
\end{reference}
Let $S$ be a scheme. Let $X = \Spec(A)$ be an affine scheme over $S$.
Let $T \subset X$ be a closed subscheme.
\begin{enumerate}
\item If the formal completion $X_{/T}$ is countably indexed
and there exist countably many $f_1, f_2, f_3, \ldots \in A$ such that
$T = V(f_1, f_2, f_3, \ldots)$, then $X_{/T}$ is adic*.
\item The conclusion of (1) is wrong if we omit the assumption that
$T$ can be cut out by countably many functions in $X$.
\end{enumerate}
\end{lemma}

\begin{proof}
The assumption that $X_{/T}$ is countably indexed means that there exists a
sequence of ideals
$$
A \supset J_1 \supset J_2 \supset J_3 \supset \ldots
$$
with $V(J_n) = T$ such that every ideal $J \subset A$ with $V(J) = T$
there exists an $n$ such that $J \supset J_n$.

\medskip\noindent
To construct an example for (2) let $\omega_1$ be the first uncountable
ordinal. Let $k$ be a field and let
$A$ be the $k$-algebra generated by $x_\alpha$, $\alpha \in \omega_1$
and $y_{\alpha \beta}$ with $\alpha \in \beta \in \omega_1$
subject to the relations $x_\alpha = y_{\alpha \beta} x_\beta$.
Let $T = V(x_\alpha)$. Let $J_n = (x_\alpha^n)$.
If $J \subset A$ is an ideal such that
$V(J) = T$, then $x_\alpha^{n_\alpha} \in J$ for some $n_\alpha \geq 1$.
One of the sets $\{\alpha \mid n_\alpha = n\}$ must be unbounded in
$\omega_1$. Then the relations imply that $J_n \subset J$.

\medskip\noindent
To see that (2) holds it now suffices to show that $A^\wedge = \lim A/J_n$
is not a ring complete with respect to a finitely generated ideal.
For $\gamma \in \omega_1$ let $A_\gamma$ be the quotient of $A$
by the ideal generated by $x_\alpha$, $\alpha \in \gamma$ and
$y_{\alpha \beta}$, $\alpha \in \gamma$. As $A/J_1$ is reduced,
every topologically nilpotent element $f$ of $\lim A/J_n$ is in
$J_1^\wedge = \lim J_1/J_n$. This means $f$ is an infinite series
involving only a countable number of generators. Hence $f$ dies in
$A_\gamma^\wedge = \lim A_\gamma/J_nA_\gamma$ for some $\gamma$.
Note that $A^\wedge \to A_\gamma^\wedge$ is continuous and open by
Lemma \href{https://stacks.math.columbia.edu/tag/0AS0}{lemma ses (Tag 0AS0)}.
If the topology on $A^\wedge$ was $I$-adic for some finitely generated ideal
$I \subset A^\wedge$, then $I$ would go to zero in some
$A_\gamma^\wedge$. This would mean that $A_\gamma^\wedge$ is discrete,
which is not the case as there is a surjective continuous and open
(by Lemma \href{https://stacks.math.columbia.edu/tag/0AS0}{lemma ses (Tag 0AS0)}) map
$A_\gamma^\wedge \to k[[t]]$ given by
$x_\alpha \mapsto t$, $y_{\alpha \beta} \mapsto 1$ for
$\gamma = \alpha$ or $\gamma \in \alpha$.

\medskip\noindent
Before we prove (1) we first prove the following: If $I \subset A^\wedge$ is
a finitely generated ideal whose closure $\bar I$ is open, then $I = \bar I$.
Since $V(J_n^2) = T$ there exists an $m$ such that $J_n^2 \supset J_m$.
Thus, we may assume that $J_n^2 \supset J_{n + 1}$ for all $n$ by passing
to a subsequence. Set $J_n^\wedge = \lim_{k \geq n} J_n/J_k \subset A^\wedge$.
Since the closure $\bar I = \bigcap (I + J_n^\wedge)$
(Lemma \href{https://stacks.math.columbia.edu/tag/0AMS}{lemma closed (Tag 0AMS)}) is open we see that there exists an $m$ such that
$I + J_n^\wedge \supset J_m^\wedge$ for all $n \geq m$. Fix such an $m$.
We have
$$
J_{n - 1}^\wedge I + J_{n + 1}^\wedge \supset
J_{n - 1}^\wedge (I + J_{n + 1}^\wedge) \supset
J_{n - 1}^\wedge J_m^\wedge
$$
for all $n \geq m + 1$. Namely, the first inclusion is trivial and the
second was shown above. Because $J_{n - 1}J_m \supset J_{n - 1}^2 \supset J_n$
these inclusions show that the image of $J_n$ in $A^\wedge$
is contained in the ideal $J_{n - 1}^\wedge I + J_{n + 1}^\wedge$.
Because this ideal is open we conclude that
$$
J_{n - 1}^\wedge I + J_{n + 1}^\wedge \supset J_n^\wedge.
$$
Say $I = (g_1, \ldots, g_t)$. Pick $f \in J_{m + 1}^\wedge$.
Using the last displayed inclusion, valid for all $n \geq m + 1$,
we can write by induction on $c \geq 0$
$$
f = \sum f_{i, c} g_i \mod J_{m + 1+ c}^\wedge
$$
with $f_{i, c} \in J_m^\wedge$ and
$f_{i, c} \equiv f_{i, c - 1} \bmod J_{m + c}^\wedge$.
It follows that $IJ_m^\wedge \supset J_{m + 1}^\wedge$.
Combined with $I + J_{m + 1}^\wedge \supset J_m^\wedge$
we conclude that $I$ is open.

\medskip\noindent
Proof of (1). Assume $T = V(f_1, f_2, f_3, \ldots)$.
Let $I_m \subset A^\wedge$ be the ideal generated by $f_1, \ldots, f_m$.
We distinguish two cases.

\medskip\noindent
Case I: For some $m$ the closure of $I_m$ is open.
Then $I_m$ is open by the result of the previous paragraph.
For any $n$ we have $(J_n)^2 \supset J_{n+1}$ by design, so
the closure of $(J_n^\wedge)^2$ contains $J_{n+1}^\wedge$
and thus is open.  Taking $n$ large, it follows that the closure
of the product of any two open ideals in $A^\wedge$ is open.
Let us prove $I_m^k$ is open for $k \ge 1$ by induction on $k$.
The case $k = 1$ is our hypothesis on $m$ in Case I.
For $k > 1$, suppose $I_m^{k - 1}$ is open. Then
$I_m^k = I_m^{k - 1} \cdot I_m$ is the product of two open ideals
and hence has open closure. But then since $I_m^k$
is finitely generated it follows that $I_m^k$
is open by the previous paragraph (applied to $I = I_m^k$),
so we can continue the induction on $k$.
As each element of $I_m$ is topologically nilpotent, we conclude
that $I_m$ is an ideal of definition which proves that $A^\wedge$
is adic with a finitely generated ideal of definition, i.e.,
$X_{/T}$ is adic*.

\medskip\noindent
Case II. For all $m$ the closure $\bar I_m$ of $I_m$ is not open.
Then the topology on $A^\wedge/\bar I_m$ is not discrete. This means
we can pick $\phi(m) \geq m$ such that
$$
\Im(J_{\phi(m)} \to A/(f_1, \ldots, f_m)) \not =
\Im(J_{\phi(m) + 1} \to A/(f_1, \ldots, f_m))
$$
To see this we have used that
$A^\wedge/(\bar I_m + J_n^\wedge) = A/((f_1, \ldots, f_m) + J_n)$.
Choose exponents $e_i > 0$ such that $f_i^{e_i} \in J_{\phi(m) + 1}$
for $0 < m < i$. Let $J = (f_1^{e_1}, f_2^{e_2}, f_3^{e_3}, \ldots)$.
Then $V(J) = T$. We claim that $J \not \supset J_n$ for all $n$
which is a contradiction proving Case II does not occur.
Namely, the image of $J$ in $A/(f_1, \ldots, f_m)$ is contained
in the image of $J_{\phi(m) + 1}$ which is properly contained in the
image of $J_m$.
\end{proof}
\end{document}
